\section{Conclusions and outlook}
\label{sec:conclusions}

This report described how \texttt{gdpr\_anonymizer} was designed, implemented, and verified as a local video-privacy toolchain. The following subsections summarise what the repository delivers today, how it behaves in practice, and which improvements remain open for future work.

\subsection{What was achieved}

The central outcome is a single public repository that accepts an MP4 file as input and produces a blurred MP4 as output, while keeping datasets, training runs, and model weights on the operator's machine and outside version control. The Python \texttt{extractor} and the Rust \texttt{segmentator} were present from the initial import of upstream tooling \cite{DecrypheExtractor,DecrypheSegmentator}; the contribution documented here is the \texttt{segmentator/ml\_pipeline} package, which closes the gap between hand-labelled polygons and an repeatable anonymisation endpoint through fourteen documented command-line stages \cite{RepoGDPR}.

At project start, the imported stack offered frame extraction, augmentation, and polygon labelling, but no in-repository training drivers, validation metrics, JSON reporting, unit-test harness, or video blur CLI. The work reported in Sections~\ref{sec:languages}--\ref{sec:testing} added dataset normalisation, portable \texttt{dataset.yaml} repair, YOLO26-seg training and validation, error analysis, optional export, explicit video blur, cleaning utilities, and warm-started iteration. Every pipeline command can emit structured JSON, which allows an operator to inspect counts, thresholds, and recommendations before deciding whether to label again, train again, or proceed to blur.

From a product perspective, the repository targets a practical privacy problem emphasised in the course presentation: sharing audiovisual material without consent exposes identifiable persons and, in many driving or street scenes, vehicle registration regions that must be obscured before publication. Regulation such as the GDPR \cite{GDPR2016} does not prescribe a specific technical method, but it reinforces that personal data should be processed lawfully and with appropriate safeguards. A documented open toolchain lowers the barrier for researchers who need reproducible control rather than ad hoc manual redaction.

To the author's knowledge, \texttt{gdpr\_anonymizer} is one of the few openly published pipelines that covers the full path from raw video through segmentation training to mask-based blur for both persons and cars in one repository. Many capable commercial products exist, yet they are often costly, closed, or limited to a single stage such as detection-only redaction or external training. The present design deliberately prioritises simplicity, explicit paths on the command line, and GitHub-hosted documentation so that others can audit, extend, or fork the workflow under the PolyForm Noncommercial licence \cite{PolyFormNC}.

Two qualitative properties distinguish the anonymisation approach from bounding-box-only redaction, as argued in the presentation material. First, instance \emph{segmentation} masks follow object contours; Gaussian or pixelated blur inside the mask union avoids the large rectangular boxes that leak background context around profile faces or partial vehicles and look conspicuous in social-media footage. Second, usable blur does not necessarily require a massive proprietary dataset at the outset. Experiments on short source clips already produced fairly convincing anonymisation when only a modest number of labelled frames were available; as more diverse footage containing pedestrians and cars is annotated, the same pipeline can yield a stronger pretrained checkpoint suitable for batch processing of new videos without redesigning the workflow.

\subsection{Does it work}

Empirical use confirms that the pipeline performs substantially better than initially expected for a student product scope. With relatively little labelled material extracted from short input videos, validation already reached bands where the interpretation layer marks \texttt{ready\_for\_video\_blur} in \texttt{metrics.json}, and the blur stage produced visually acceptable output on multi-thousand-frame clips. Reported blur KPIs include frame counts, detection rates, per-class totals, and average masked area fraction in \texttt{runs/blur/report.json}, which makes downstream review straightforward.

Quality still scales with data and labelling effort, as in any supervised segmentation project. More annotated frames, broader scene diversity, and careful polygon boundaries improve mask mAP and reduce missed or imprecise regions. The repository supports that growth without forcing the operator to restart from scratch: \texttt{gdpr-yolo-iterate} warm-starts a new training run from an existing \texttt{best.pt}, re-validates on the same split, and records metric deltas against a prior validation report. Combined with \texttt{gdpr-yolo-analyze-errors}, which surfaces worst validation images, the workflow implements a feedback loop from metrics back to labelling and training rather than a one-shot train-and-hope procedure.

Deployment remains local by design. Training and blur require the optional Ultralytics and PyTorch stack, yet neither a dedicated cloud service nor high-end GPU hardware is mandatory. CPU execution is supported throughout; a GPU mainly reduces wall-clock time on long videos or large epoch counts. Continuous integration on GitHub Actions runs the Rust and Python unit suites on every pull request \cite{CIGDPR}, and the coverage evolution summarised in Section~\ref{sec:testing} indicates that regressions in CLI drivers and video logic are unlikely to pass unnoticed.

\subsection{What remains to be done}

Several limitations are understood and documented rather than hidden. The most visible artefact in plate and small-object blur is temporal \emph{flickering}: when segmentation masks shift abruptly between neighbouring frames, the obscured region appears to jump even though each frame is processed correctly in isolation. A practical mitigation---not yet implemented---is temporal propagation of stable polygon or mask geometry. If frame~3 and frames~7--9 share substantially the same segment coordinates for a given instance, intermediate frames can inherit interpolated or held positions so that blur boundaries remain consistent across the clip. This is compatible with the existing per-frame blur architecture and would not require abandoning the current JSON-first CLI design.

Further gains depend primarily on dataset curation, which remains the responsibility of each operator. Better masks require more accurate polygons, especially on profile faces, partially occluded plates, and varied lighting. Presentation notes highlighted additional data axes: more pedestrian diversity, more licence-plate examples, plates from different countries, and enough epochs to stabilise small moving regions. Those are process recommendations rather than missing code paths; the augmentor, segmentator, and iteration tools already support incremental dataset growth.

Tooling inside the segmentator could also reduce labelling friction. A ``magic wand'' or assisted edge selection around faces and plate borders would accelerate polygon creation and improve boundary accuracy compared with purely manual vertex placement. Keyboard shortcuts, zoom, and scroll-bar improvements were already added to the Rust UI during the project; assisted selection would be a natural next ergonomic step.

Finally, integrating external detection datasets that ship bounding boxes but not segmentation polygons would require a conversion stage, as noted when evaluating public Roboflow face-detection corpora during the presentation. Such a step would sit beside the existing fourteen-stage table rather than replace mask-based training, because box-only labels cannot reproduce the contour-following blur that motivates the segmentation model choice.

In summary, \texttt{gdpr\_anonymizer} delivers a coherent, test-covered, locally runnable path from video ingest to GDPR-oriented anonymisation output. It already demonstrates that modest labelled data can produce useful blur, that segmentation masks outperform coarse boxes for visual quality, and that iteration hooks prevent dead-end training cycles. Remaining work centres on temporal stability in video blur, richer labelled data, and faster annotation assistance---extensions that build on the foundation described in this report rather than replacing it.
